\section{Discussion}
\label{sec:discussion}

\subsection{A Functionally Decoupled View of Ambiguity}

The experimental results support the central premise of this work: ambiguity in visual intention recognition is not a single phenomenon, but a functionally decoupled one. Supervisory ambiguity destabilizes the semantic supervision signal itself, whereas semantic ambiguity distorts the local decision geometry among neighboring intent labels. \revised{This is why the manner of combining the two signals matters. The semantic prior and the rationale teacher are each only partial on their own, and forcing both into a single undifferentiated supervision target underperforms the decoupled model. The benefit of separating the two functions lies less in a higher accuracy ceiling than in keeping each function individually justified, separately measurable, and independently deployable.}

This viewpoint also clarifies why a strong CLIP baseline already performs well. Once visual representation is sufficiently expressive, the remaining bottleneck is no longer basic visual recognition, but the mapping from rich visual semantics to abstract intent labels. In other words, the main challenge is not merely seeing better, but deciding better under ambiguity.

\subsection{\revised{Main Findings and Advantages}}

\revised{Our experiments support four main findings. First, the gains are genuine rather than an artifact of calibration or backbone strength: under a matched CLIP ViT-L/14 backbone and an identical validation-only thresholding protocol, FDIL improves the CLIP baseline by $+4.1$ average F1 and $+6.6$ hard-subset F1, with the threshold-free mAP improving by $+4.4$ and a bootstrap confidence interval that excludes zero. Second, the two modules are empirically separable, each contributing an independent increment on threshold-free mAP ($+2.4$ for UTD and $+2.0$ for SLR-C) that mirrors the two-term risk decomposition of Section~\ref{sec:decoupling_principle}. Third, the benefit comes from rationale content rather than from external text per se: replacing each rationale with a mismatched one is the single most damaging intervention we tried, worse than removing the teacher altogether. Fourth, the improvements concentrate where the task is hardest, with the largest per-class gains on abstract, disagreement-prone intents and only two of twenty-eight classes regressing. The practical advantages that follow are a clean two-function design, a UTD branch deployable alone at zero extra inference cost, and a full model whose feature-level overhead stays well below one millisecond per image.}

% \subsection{Heterogeneous Priors Improve Robustness Only with Controlled Integration}

% Our results show that heterogeneous priors are useful, but not in an unconstrained way. In SLR-C, the lexical, canonical, and scenario sources provide complementary signals: lexical phrases are concise but coarse, canonical descriptions are semantically cleaner but still abstract, and scenario descriptions offer the strongest visual grounding. Their combination improves semantic coverage and hard-case robustness. However, the component ablation also shows that a fixed heterogeneous prior alone is insufficient and may even distort the score geometry when directly injected without adaptive refinement.

% This observation is important for understanding the role of SLR-C in the final FDIL system. The prior should not be interpreted as a standalone classifier. Instead, it should be viewed as a structured semantic organizer that proposes a more informed local score space, after which residual refinement adapts the prediction to image evidence. This is precisely why the full system outperforms naive prior fusion, and why heterogeneous semantics become most effective when coupled with trainable adaptation.

% \subsection{Text Teacher Distillation Provides Structured Semantic Supervision}

% The rationale-related ablations reveal another important point: not every form of text supervision is equally useful. Generic image captions already provide some privileged signal, but they remain weaker than structured rationale. Likewise, simply enriching the text with contextual interpretation is not enough; the largest gain appears only when counterfactual disambiguation is included. This suggests that the most useful textual signal for subjective intent recognition is not raw description, but a semantically sharpened explanation that explicitly contracts the boundary between nearby intents.

% The modality ablation further reinforces this interpretation. A multimodal teacher with both image and text inputs is not the strongest distillation target, even though its oracle performance is not necessarily worse. This implies that the teacher is most useful when it behaves as a cleaner semantic supervisor rather than as another fused classifier. In this sense, UTD works less like conventional multimodal fusion and more like structured semantic regularization.

\subsection{Limitations and Future Directions}

Despite the consistent gains, the current framework still has several limitations. First, the method depends on offline text generation and external prior construction, which increases preparation cost even though UTD itself introduces no extra inference-time cost. Second, the gains shrink on weaker backbones such as ResNet101, suggesting that the current residual refinement strategy still relies on reasonably strong base representations. Third, while the three-source prior improves robustness, its behavior also shows that richer priors are not automatically better unless their interaction with the student is carefully controlled. \revised{Fourth, the per-class attribution reveals a characteristic failure mode: for a few broadly defined, mutually overlapping categories (most visibly \textit{CreativeUnique}), the residual student can resolve the local candidate set toward a competing neighbor and lose accuracy on the original class, so local reranking does not fully disentangle every tightly coupled confusion group. Fifth, the decomposition is validated on visual intent and transfers only partially to the EMOTIC affect dataset. Because EMOTIC multi-annotates only its validation and test splits, the single-annotator training labels make per-sample agreement degenerate on the standard split and the agreement-gated term reduces to standard distillation by construction. In a controlled cross-validation on EMOTIC's multi-annotated val+test pool, where per-sample agreement is available at training time, privileged distillation significantly improves the image-only baseline on all metrics (paired $p<0.01$), evidencing framework-level transfer to a second domain; the gate, however, only matches standard distillation there, where the borrowed agreement is coarser than Intentonomy's native training-time agreement. We therefore claim framework-level generalization while reserving the gate's incremental benefit for settings with rich training-time agreement, and do not claim that FDIL transfers automatically to every subjective multi-label task.}

These limitations suggest several future directions. One is to learn source-aware weighting instead of uniform prior averaging, so that the model can adaptively choose which semantic source to trust for each class or sample. Another is to design stronger student architectures for weaker backbones, making the functionally decoupled framework less dependent on CLIP-level representations. The third is to incorporate more explicit candidate-level reasoning, so that textual guidance operates not only through global distillation and local reranking, but also through confusion-pair-specific evidence resolution.
